\documentclass[12pt,a4paper]{article}
\usepackage[utf8]{inputenc}
\usepackage[T2A]{fontenc}
\usepackage[russian]{babel}
\usepackage{amsmath, amssymb, amsthm}
\usepackage{mathtools}
\usepackage{array}
\usepackage{booktabs}
\usepackage[hidelinks]{hyperref}
\usepackage{geometry}
\geometry{margin=2.5cm}

\newtheorem{definition}{Определение}
\newtheorem{predicate}[definition]{Предикат}

\title{%
  Алгоритм DST-агрегации для анализа компетентностных дефицитов:\\
  \large Формальное описание в предикатной логике и теории множеств
}
\author{}
\date{}

% ─────────────────────────────────────────────────────────────────────────────
\begin{document}
\maketitle
\tableofcontents
\bigskip

% =============================================================================
\section{Базовые множества и рамка рассуждений}
% =============================================================================

\begin{definition}[Рамка рассуждений]
  Для одного кластера компетенций рамка рассуждений определяется как
  \[
    \Theta = \{T,\, F\},
  \]
  где $T$ означает «кластер важен» (True), $F$ --- «кластер не важен» (False).
  Степенное множество $2^\Theta = \{\emptyset,\, \{T\},\, \{F\},\, \Theta\}$.
\end{definition}

\begin{definition}[Множество источников]
  \[
    \mathcal{S} = \{\text{vac},\; \text{exp},\; \text{fc}\},
  \]
  где \textbf{vac} --- данные вакансий, \textbf{exp} --- экспертные оценки,
  \textbf{fc} --- прогнозы (forecasts).
\end{definition}

\begin{definition}[Входной кортеж источника]
  Каждый источник $s \in \mathcal{S}$ задаётся кортежем
  \[
    \mathcal{D}_s
      = \bigl(n_s^{\mathrm{tot}},\; n_s^{\mathrm{rel}},\;
              \bar{a}_s,\; \alpha_s,\; w_s,\; \lambda_s\bigr),
  \]
  где:
  \begin{itemize}
    \item $n_s^{\mathrm{tot}} \in \mathbb{N}$ --- общее число записей источника,
    \item $n_s^{\mathrm{rel}} \in \mathbb{N},\ n_s^{\mathrm{rel}} \le n_s^{\mathrm{tot}}$ --- записей, относящихся к кластеру,
    \item $\bar{a}_s \in [0,1]$ --- средняя оценка важности кластера,
    \item $\alpha_s \in [0,1]$ --- уровень согласованности источника,
    \item $w_s \in [0,1]$ --- вес надёжности источника,
    \item $\lambda_s > 0$ --- константа сглаживания.
  \end{itemize}

  Для прогнозов со смешанным направлением вводится расширенный кортеж
  \[
    \mathcal{D}_{\mathrm{fc}}^{\mathrm{mix}}
      = \bigl(n_{\mathrm{fc}}^{\mathrm{tot}},\;
              n_+,\; \bar{a}_+,\;
              n_-,\; \bar{a}_-,\;
              w_{\mathrm{fc}},\; \lambda_{\mathrm{fc}}\bigr),
  \]
  где $n_+$ ($n_-$) --- количество прогнозов роста (спада),
  $\bar{a}_+$ ($\bar{a}_-$) --- средние оценки соответственно,
  $n_+ + n_- = n_{\mathrm{fc}}^{\mathrm{tot}}$.

\end{definition}

\begin{definition}[Обеспечённость кластера]
  $\sigma \in [0,1]$ --- текущая доля кластера в учебной программе
  (РПД / учебный план).
\end{definition}

% =============================================================================
\section{Шаг 1. Коэффициент уверенности $\kappa$}
% =============================================================================

\subsection{Мотивация}

Источник с тысячами записей заслуживает большего доверия, чем источник
с тремя. При этом уверенность ограничена сверху: $\kappa \in (0,1)$.

\subsection{Формулы}

\begin{definition}[Эффективный объём источника]
  \[
    n_s^{\mathrm{eff}} = n_s^{\mathrm{rel}} \cdot \alpha_s.
  \]
  Для смешанных прогнозов ($\alpha_{\mathrm{fc}} = 1$):
  \[
    n_{\mathrm{fc}}^{\mathrm{eff}} = n_+ + n_-.
  \]
\end{definition}

\begin{definition}[Коэффициент уверенности]
  \[
    \kappa_s = \frac{n_s^{\mathrm{eff}}}{n_s^{\mathrm{eff}} + \lambda_s}.
  \]
\end{definition}

\subsection{Свойства}

\begin{align*}
  n_s^{\mathrm{eff}} = 0     &\;\Rightarrow\; \kappa_s = 0
    && \text{(нет данных --- нет уверенности)}, \\
  n_s^{\mathrm{eff}} \to \infty &\;\Rightarrow\; \kappa_s \to 1
    && \text{(бесконечный объём --- полная уверенность)}, \\
  n_s^{\mathrm{eff}} = \lambda_s &\;\Rightarrow\; \kappa_s = 0{,}5
    && \text{(«полудоверие» --- смысл $\lambda_s$)}.
\end{align*}

\subsection{Зафиксированные значения $\lambda_s$}

\begin{center}
\begin{tabular}{lcc}
\toprule
Источник $s$ & $\lambda_s$ & Смысл \\
\midrule
vac & 15 & 15 вакансий $\Rightarrow \kappa = 0{,}5$ \\
exp & 5  & 5 экспертов $\Rightarrow \kappa = 0{,}5$ \\
fc  & 2  & 2 прогноза $\Rightarrow \kappa = 0{,}5$ \\
\bottomrule
\end{tabular}
\end{center}

% =============================================================================
\section{Шаг 2. Базовое распределение вероятности (BPA)}
% =============================================================================

\subsection{Определение}

\begin{definition}[BPA]
  Функция базового распределения вероятности для источника $s$:
  \[
    m_s : 2^\Theta \to [0,1],
    \qquad
    \sum_{A \subseteq \Theta} m_s(A) = 1,
    \qquad
    m_s(\emptyset) = 0.
  \]
\end{definition}

\subsection{Единая схема BPA для всех источников}

Для каждого источника $s$ всегда вычисляются три составляющих:

\begin{align}
  m_s\!\bigl(\{T\}\bigr)
    &= \kappa_s \cdot \bar{a}_s^{+} \cdot
       \frac{n_s^{+}}{n_s^{\mathrm{tot}}}, \label{eq:mT} \\[6pt]
  m_s\!\bigl(\{F\}\bigr)
    &= \kappa_s \cdot \bar{a}_s^{-} \cdot
       \frac{n_s^{-}}{n_s^{\mathrm{tot}}}, \label{eq:mF} \\[6pt]
  m_s(\Theta)
    &= 1 - m_s\!\bigl(\{T\}\bigr) - m_s\!\bigl(\{F\}\bigr), \label{eq:mU}
\end{align}

где $n_s^{+}$, $\bar{a}_s^{+}$ --- количество и средняя оценка
\emph{положительных} сигналов источника $s$;
$n_s^{-}$, $\bar{a}_s^{-}$ --- отрицательных.

\paragraph{Частный случай для vac и exp.}
Данные источники содержат только подтверждающие сигналы:
$n_s^{-} = 0$, поэтому $m_s(\{F\}) = 0$ (не аксиома, а следствие данных).
Формулы \eqref{eq:mT}--\eqref{eq:mU} остаются теми же.

\paragraph{Частный случай для fc (смешанный тип).}
Прогнозы роста дают $n_{\mathrm{fc}}^{+} > 0$, прогнозы спада ---
$n_{\mathrm{fc}}^{-} > 0$, поэтому оба слагаемых $m_{\mathrm{fc}}(\{T\})$
и $m_{\mathrm{fc}}(\{F\})$ ненулевые.

\subsection{Предикат корректности BPA}

\begin{predicate}[Нормировка]
  \[
    \forall s \in \mathcal{S}:\quad
    m_s\!\bigl(\{T\}\bigr) + m_s\!\bigl(\{F\}\bigr) + m_s(\Theta) = 1
    \;\wedge\;
    m_s(\emptyset) = 0.
  \]
\end{predicate}

% =============================================================================
\section{Шаг 3. Дисконтирование}
% =============================================================================

\subsection{Определение}

Дисконтирование «перекладывает» часть уверенности источника
в неопределённость пропорционально его ненадёжности $(1 - w_s)$.

\begin{definition}[Дисконтированное BPA]
  \begin{align}
    m'_s(A)      &= w_s \cdot m_s(A),
      \quad \forall A \subsetneq \Theta,
      \label{eq:disc_A} \\[6pt]
    m'_s(\Theta) &= 1 - w_s \cdot \bigl(1 - m_s(\Theta)\bigr)
                 = (1 - w_s) + w_s \cdot m_s(\Theta).
      \label{eq:disc_Theta}
  \end{align}
\end{definition}

\noindent
Для $\Theta = \{T, F\}$ раскрываем явно:
\begin{align*}
  m'_s\!\bigl(\{T\}\bigr) &= w_s \cdot m_s\!\bigl(\{T\}\bigr), \\
  m'_s\!\bigl(\{F\}\bigr) &= w_s \cdot m_s\!\bigl(\{F\}\bigr), \\
  m'_s(\Theta)             &= (1 - w_s) + w_s \cdot m_s(\Theta).
\end{align*}

\subsection{Зафиксированные веса $w_s$}

\begin{center}
\begin{tabular}{lcc}
\toprule
Источник $s$ & $w_s$ & Обоснование \\
\midrule
exp  & $0{,}90$ & Профессиональные оценки, высокая надёжность \\
vac  & $0{,}80$ & Фактические данные, есть дубли и сезонность \\
fc   & $0{,}60$ & Прогнозы часто ошибаются в IT-отрасли \\
\bottomrule
\end{tabular}
\end{center}

% =============================================================================
\section{Шаг 4. Комбинирование по правилу Демпстера}
% =============================================================================

\subsection{Коэффициент конфликта}

\begin{definition}[Коэффициент конфликта двух BPA]
  Для $m_1, m_2 : 2^\Theta \to [0,1]$:
  \[
    K(m_1, m_2)
      = \sum_{\substack{B,\, C \,\subseteq\, \Theta \\ B \,\cap\, C \,=\, \emptyset}}
          m_1(B) \cdot m_2(C).
  \]
  Для $\Theta = \{T, F\}$ это раскрывается в:
  \[
    K(m_1, m_2)
      = m_1\!\bigl(\{T\}\bigr) \cdot m_2\!\bigl(\{F\}\bigr)
      + m_1\!\bigl(\{F\}\bigr) \cdot m_2\!\bigl(\{T\}\bigr).
  \]
\end{definition}

\subsection{Предикат низкого конфликта}

\begin{predicate}[Допустимый конфликт]
  \[
    \mathrm{LowConflict}(m_1, m_2)
      \;\Longleftrightarrow\;
      K(m_1, m_2) < \tau_K.
  \]
\end{predicate}

\subsection{Правило Демпстера (при низком конфликте)}

\begin{definition}[Комбинирование по Демпстеру]
  Если $\mathrm{LowConflict}(m_1, m_2)$:
  \[
    (m_1 \oplus m_2)(A)
      = \frac{1}{1 - K(m_1, m_2)}
        \sum_{\substack{B \,\cap\, C \,=\, A \\ A \neq \emptyset}}
          m_1(B) \cdot m_2(C),
    \qquad A \neq \emptyset.
  \]
\end{definition}

Для $\Theta = \{T, F\}$ раскрываем:

\begin{align}
  (m_1 \oplus m_2)\!\bigl(\{T\}\bigr)
    &= \frac{
        m_1(\{T\}) \cdot m_2(\{T\})
        + m_1(\{T\}) \cdot m_2(\Theta)
        + m_1(\Theta) \cdot m_2(\{T\})
      }{1 - K},
    \label{eq:comb_T} \\[8pt]
  (m_1 \oplus m_2)(\Theta)
    &= \frac{m_1(\Theta) \cdot m_2(\Theta)}{1 - K},
    \label{eq:comb_U} \\[8pt]
  (m_1 \oplus m_2)\!\bigl(\{F\}\bigr)
    &= \frac{
        m_1(\{F\}) \cdot m_2(\{F\})
        + m_1(\{F\}) \cdot m_2(\Theta)
        + m_1(\Theta) \cdot m_2(\{F\})
      }{1 - K}.
    \label{eq:comb_F}
\end{align}

\subsection{Последовательное комбинирование трёх источников}

\[
  m^{(1)} = m'_{\mathrm{vac}} \oplus m'_{\mathrm{exp}},
  \qquad
  K_1 = K\!\bigl(m'_{\mathrm{vac}},\, m'_{\mathrm{exp}}\bigr),
\]
\[
  m^* = m^{(1)} \oplus m'_{\mathrm{fc}},
  \qquad
  K_2 = K\!\bigl(m^{(1)},\, m'_{\mathrm{fc}}\bigr),
\]
\[
  K = \max(K_1,\, K_2).
\]

Правило Демпстера коммутативно и ассоциативно:
\[
  m'_{\mathrm{vac}} \oplus m'_{\mathrm{exp}} \oplus m'_{\mathrm{fc}}
  = m'_{\mathrm{exp}} \oplus m'_{\mathrm{vac}} \oplus m'_{\mathrm{fc}}
  = \cdots
\]

% =============================================================================
\section{Шаг 5. Пигнистическая вероятность и показатель дефицита}
% =============================================================================

\subsection{BetP}

\begin{definition}[Pignistic Probability]
  В рамке $\Theta = \{T, F\}$ из $N_C$ кластеров программы:
  \[
    \mathrm{BetP}(T)
      = m^*\!\bigl(\{T\}\bigr)
        + \frac{m^*(\Theta)}{N_C},
  \]
  где $N_C$ --- количество кластеров в учебной программе
  (текущее значение: $N_C = 25$).
\end{definition}

\noindent
\textit{Смысл:} неопределённая масса $m^*(\Theta)$ равномерно
распределяется между всеми $N_C$ кластерами.
Общая формула через суммирование по фокальным элементам:
\[
  \mathrm{BetP}(T)
    = \sum_{\substack{A \,\subseteq\, \Theta \\ T \,\in\, A}}
        \frac{m^*(A)}{|A|}
    = m^*\!\bigl(\{T\}\bigr) + \frac{m^*(\Theta)}{|\Theta|}.
\]

\subsection{Показатель дефицита}

\begin{definition}[Дефицит кластера]
  \[
    \Delta = \mathrm{BetP}(T) - \sigma.
  \]
\end{definition}

\subsection{Интерпретация}

\[
  \Delta > \tau_\Delta
    \;\Rightarrow\; \text{дефицит: рынок требует больше, чем есть в РПД},
\]
\[
  |\Delta| \le \tau_\Delta
    \;\Rightarrow\; \text{баланс},
\]
\[
  \Delta < -\tau_\Delta
    \;\Rightarrow\; \text{избыток: в РПД больше, чем требует рынок}.
\]

% =============================================================================
\section{Шаг 6. Правила принятия решения}
% =============================================================================

\subsection{Три условия-предиката}

\begin{predicate}[Дефицит]
  \[
    \mathrm{Def} \;\Longleftrightarrow\; \Delta > \tau_\Delta.
  \]
\end{predicate}

\begin{predicate}[Согласованность источников]
  \[
    \mathrm{Cons} \;\Longleftrightarrow\; K \le \tau_K.
  \]
\end{predicate}

\begin{predicate}[Достаточная определённость]
  \[
    \mathrm{Cert} \;\Longleftrightarrow\; m^*(\Theta) \le \tau_\Theta.
  \]
\end{predicate}

\subsection{Таблица решений}

\begin{center}
\renewcommand{\arraystretch}{1.4}
\begin{tabular}{ccc l}
\toprule
$\mathrm{Def}$ & $\mathrm{Cons}$ & $\mathrm{Cert}$ & Решение \\
\midrule
$\top$ & $\top$ & $\top$ & \textbf{УСИЛИТЬ} \\
$\top$ & $\top$ & $\bot$ & \textbf{ВАРИАТИВНОСТЬ} (высокая неопределённость) \\
$\top$ & $\bot$ & любое  & \textbf{ВАРИАТИВНОСТЬ} (конфликт источников) \\
$\bot$ ($\Delta < -\tau_\Delta$) & --- & --- & \textbf{СОКРАТИТЬ} \\
$\bot$ ($|\Delta| \le \tau_\Delta$) & --- & --- & \textbf{СТАБИЛИЗИРОВАТЬ} \\
\bottomrule
\end{tabular}
\end{center}

\subsection{Формальные правила в предикатной логике}

\begin{align}
  \mathrm{REINFORCE}
    &\;\Longleftrightarrow\;
      \Delta > \tau_\Delta
      \;\wedge\; K \le \tau_K
      \;\wedge\; m^*(\Theta) \le \tau_\Theta,
    \label{eq:reinforce} \\[6pt]
  \mathrm{VARIABILITY}
    &\;\Longleftrightarrow\;
      \Delta > \tau_\Delta
      \;\wedge\; \bigl(K > \tau_K
                       \;\vee\; m^*(\Theta) > \tau_\Theta\bigr),
    \label{eq:variability} \\[6pt]
  \mathrm{REDUCE}
    &\;\Longleftrightarrow\;
      \Delta < -\tau_\Delta,
    \label{eq:reduce} \\[6pt]
  \mathrm{STABILIZE}
    &\;\Longleftrightarrow\;
      |\Delta| \le \tau_\Delta,
    \label{eq:stabilize}
\end{align}

\noindent
при этом $\mathrm{REINFORCE} \cup \mathrm{VARIABILITY}$ --- непересекающееся разбиение
для случая $\Delta > \tau_\Delta$.

\subsection{Градация УСИЛИТЬ}

\[
  \mathrm{STRONG} \;\Longleftrightarrow\; \mathrm{REINFORCE} \;\wedge\; \Delta > 0{,}5,
\]
\[
  \mathrm{MODERATE} \;\Longleftrightarrow\; \mathrm{REINFORCE} \;\wedge\;
    \tau_\Delta < \Delta \le 0{,}5.
\]

% =============================================================================
\section{Сводная цепочка алгоритма}
% =============================================================================

\noindent\textbf{Вход:} $\mathcal{D}_s = (n_s^{\mathrm{tot}},\; n_s^{+},\; \bar{a}_s^{+},\; n_s^{-},\; \bar{a}_s^{-},\; \alpha_s,\; w_s,\; \lambda_s)$, $\sigma$

\begin{align*}
  (1)\quad & n_s^{\mathrm{eff}} = n_s^{\mathrm{rel}} \cdot \alpha_s,
             \qquad
             \kappa_s = \frac{n_s^{\mathrm{eff}}}{n_s^{\mathrm{eff}} + \lambda_s} \\[6pt]
  (2)\quad & m_s(\{T\}) = \kappa_s \cdot \bar{a}_s^{+} \cdot
               \frac{n_s^{+}}{n_s^{\mathrm{tot}}},\quad
             m_s(\{F\}) = \kappa_s \cdot \bar{a}_s^{-} \cdot
               \frac{n_s^{-}}{n_s^{\mathrm{tot}}},\quad
             m_s(\Theta) = 1 - m_s(\{T\}) - m_s(\{F\}) \\[6pt]
  (3)\quad & m'_s(\{T\}) = w_s \cdot m_s(\{T\}),\quad
             m'_s(\Theta) = (1-w_s) + w_s \cdot m_s(\Theta) \\[6pt]
  (4)\quad & m^* = m'_{\mathrm{vac}} \oplus m'_{\mathrm{exp}} \oplus m'_{\mathrm{fc}},
             \qquad K = \max(K_1,\, K_2) \\[6pt]
  (5)\quad & \mathrm{BetP} = m^*(\{T\}) + \frac{m^*(\Theta)}{N_C},\quad
             \Delta = \mathrm{BetP} - \sigma
\end{align*}

\noindent\textbf{Выход:} решение по формулам \eqref{eq:reinforce}--\eqref{eq:stabilize}.

% =============================================================================
\section{Параметры алгоритма}
% =============================================================================

\begin{center}
\renewcommand{\arraystretch}{1.3}
\begin{tabular}{llll}
\toprule
Параметр & Значение & Источник & Рекомендуемый метод \\
\midrule
$\lambda_{\mathrm{vac}}$ & 15 & хардкод & Медиана объёмов вакансий \\
$\lambda_{\mathrm{exp}}$ & 5  & хардкод & Медиана числа экспертов \\
$\lambda_{\mathrm{fc}}$  & 2  & хардкод & Медиана числа прогнозов \\
$w_{\mathrm{exp}}$  & $0{,}90$ & хардкод & Калибровка на истории \\
$w_{\mathrm{vac}}$  & $0{,}80$ & хардкод & Калибровка на истории \\
$w_{\mathrm{fc}}$   & $0{,}60$ & хардкод & Калибровка на истории \\
$\tau_\Delta$       & $0{,}15$ & статья  & Калибровка \\
$\tau_K$            & $0{,}40$ & статья  & Калибровка \\
$\tau_\Theta$       & $0{,}15$ & статья  & Калибровка \\
$N_C$               & 25       & хардкод & Из БД (число кластеров программы) \\
\bottomrule
\end{tabular}
\end{center}

% =============================================================================
\end{document}
